\section{Claim-by-Claim Integrated Argument}
\label{sec:oc133-claim-by-claim-integrated-argument}
This section is the prose counterpart to the claim and theorem registers. It is intentionally longer than a manifest because a publication-grade monograph must teach the reader why each promoted claim exists, what problem it solves, which evidence carries it, and where the claim stops. The register remains available for audit; the argument below is written for scientific reading.

\subsection{Shared Claim Reading Method}
Every promoted model-core claim is read by the same method. First, identify the scientific ambiguity that the claim removes. Second, identify the typed OC construction that carries the distinction. Third, identify the proof sheet, formal fragment, finite semantic witness, replay row, comparator row, or claim-boundary record that can carry the public sentence. Fourth, identify the reopening condition. This common method is stated once here so that the following subsections can focus on the scientific work of each claim rather than repeat the same editorial apparatus.

The dependency route is layered rather than interchangeable. Definitions provide vocabulary; proof sheets state assumptions and lemmas; Lean records the encoded fragment; finite semantic cases show executable positive and negative behavior where the theorem is finite; empirical replay rows support only bounded artifact and reconstruction claims; comparator rows protect novelty wording; and reviewer gates check that the public sentence has not drifted beyond the evidence. No layer is allowed to impersonate another.

The hostile-review method is likewise fixed. A reviewer may remove a tuple component, collapse a status distinction, turn a boundary into a metaphor, treat a finite witness as a broad empirical prediction, or upgrade a comparator overlap into priority. A claim survives only while the named proof and evidence path blocks that attack at the stated strength. Otherwise the claim is repaired, demoted, or held for a later research release.

The editorial method is different from the proof method but it is not optional. The prose must let a reader see the problem, construction, evidence, and boundary without parsing a control table. Row identifiers serve as citations for audit and reproduction; they are not the voice of the scientific argument. The following subsections therefore preserve exact evidence identifiers while making each claim readable as science.

\subsection{Lifecycle status: liveness, death, residue, and rebirth}
The first question for this claim is what weaker vocabulary would confuse. The claim addresses this ambiguity: Death blocks live status; residue and rebirth are token-bound evidence relations with distinct class-specific endpoint rules: residue separates source from residue while returning to the source endpoint, and rebirth separates source, residue, and new target tokens. Rebirth is non-identity unless endpoint-bound identity evidence has identity class, declared invariant preservation, no residue token, and equal source/target endpoint evidence. No categorical Hom/composition theorem is promoted. The point of including it in the 1.3.3 release is to make the ambiguity auditable rather than to decorate the theory with another label.

The support class is Lean subset, structured proof sheet, and finite witness. The proof route is public proof sheet for the lifecycle status theorem; the formal route is Lean declaration lifecycle\_residue\_rebirth\_morphism\_boundary. In the public reading, this means that the claim is carried by a named theorem boundary, not by the fact that a row exists. If the proof sheet, Lean reference, finite witness, or mutation control does not carry the visible wording, the visible wording is the part that must move.

The claim boundary is Bounded to stated theorem assumptions, finite witnesses, and public falsifier boundary.. This boundary is not decorative. It prevents a bounded model-core statement from becoming a claim about unrestricted totality, full-scope prediction, or global priority. The scientific value of the claim is that it is strong enough to test a real structural distinction and narrow enough that an adversarial reviewer can name the condition that would reopen it.

A worked reading asks three local questions. Which typed field, operator, boundary, status predicate, or morphism is doing the work? Which situation is the positive case? Which negative case would have collapsed the distinction if the claim were only a slogan? The answer to those questions is the practical content of this subsection.

For T133-OMEGA-STATUS, the corresponding falsifier is concrete rather than slogan-like: find a model satisfying the stated assumptions in which the asserted distinction collapses, show that the positive witness was accepted for the wrong reason, show that the negative control is non-responsive, or show that the public wording requires an assumption absent from the proof sheet.

The prior-art and empirical boundaries are local to the claim. Earlier formalisms may already carry part of the distinction; when they do, OC claims integration only for that part. Numeric replay rows can demonstrate artifact discipline, comparator behavior, and falsifier plumbing, but they do not by themselves complete a domain science. The reader should therefore leave with four handles: the conceptual problem, the typed construction, the evidence route, and the reopening route.

\subsection{Continuumness-zero obstruction}
The claim is useful only if it separates two situations that a less typed account would merge. The claim addresses this ambiguity: Continuumness zero requires live support, an independently clear obstruction register, and a declared zero-cause family; a zero-cause label alone does not compute k=0. The point of including it in the 1.3.3 release is to make the ambiguity auditable rather than to decorate the theory with another label.

The support class is Lean subset, structured proof sheet, and finite witness. The proof route is public proof sheet for the K-zero boundary theorem; the formal route is Lean declaration continuumness\_zero\_case\_iff\_declared\_zero\_cause\_with\_support. In the public reading, this means that the claim is carried by a named theorem boundary, not by the fact that a row exists. If the proof sheet, Lean reference, finite witness, or mutation control does not carry the visible wording, the visible wording is the part that must move.

The claim boundary is Bounded to stated theorem assumptions, finite witnesses, and public falsifier boundary.. This boundary is not decorative. It prevents a bounded model-core statement from becoming a claim about unrestricted totality, full-scope prediction, or global priority. The scientific value of the claim is that it is strong enough to test a real structural distinction and narrow enough that an adversarial reviewer can name the condition that would reopen it.

A worked reading asks three local questions. Which typed field, operator, boundary, status predicate, or morphism is doing the work? Which situation is the positive case? Which negative case would have collapsed the distinction if the claim were only a slogan? The answer to those questions is the practical content of this subsection.

For T133-K-ZERO, the corresponding falsifier is concrete rather than slogan-like: find a model satisfying the stated assumptions in which the asserted distinction collapses, show that the positive witness was accepted for the wrong reason, show that the negative control is non-responsive, or show that the public wording requires an assumption absent from the proof sheet.

The prior-art and empirical boundaries are local to the claim. Earlier formalisms may already carry part of the distinction; when they do, OC claims integration only for that part. Numeric replay rows can demonstrate artifact discipline, comparator behavior, and falsifier plumbing, but they do not by themselves complete a domain science. The reader should therefore leave with four handles: the conceptual problem, the typed construction, the evidence route, and the reopening route.

\subsection{Metric-threshold boundary specialization}
The scientific burden here is to show why the named distinction is not merely verbal. The claim addresses this ambiguity: Metric thresholds are a specialization of typed classifier boundaries. The point of including it in the 1.3.3 release is to make the ambiguity auditable rather than to decorate the theory with another label.

The support class is Lean subset, structured proof sheet, and finite witness. The proof route is public proof sheet for the boundary representation theorem; the formal route is Lean declaration metric\_boundary\_specialization. In the public reading, this means that the claim is carried by a named theorem boundary, not by the fact that a row exists. If the proof sheet, Lean reference, finite witness, or mutation control does not carry the visible wording, the visible wording is the part that must move.

The claim boundary is Bounded to stated theorem assumptions, finite witnesses, and public falsifier boundary.. This boundary is not decorative. It prevents a bounded model-core statement from becoming a claim about unrestricted totality, full-scope prediction, or global priority. The scientific value of the claim is that it is strong enough to test a real structural distinction and narrow enough that an adversarial reviewer can name the condition that would reopen it.

A worked reading asks three local questions. Which typed field, operator, boundary, status predicate, or morphism is doing the work? Which situation is the positive case? Which negative case would have collapsed the distinction if the claim were only a slogan? The answer to those questions is the practical content of this subsection.

For T133-BOUNDARY, the corresponding falsifier is concrete rather than slogan-like: find a model satisfying the stated assumptions in which the asserted distinction collapses, show that the positive witness was accepted for the wrong reason, show that the negative control is non-responsive, or show that the public wording requires an assumption absent from the proof sheet.

The prior-art and empirical boundaries are local to the claim. Earlier formalisms may already carry part of the distinction; when they do, OC claims integration only for that part. Numeric replay rows can demonstrate artifact discipline, comparator behavior, and falsifier plumbing, but they do not by themselves complete a domain science. The reader should therefore leave with four handles: the conceptual problem, the typed construction, the evidence route, and the reopening route.

\subsection{Typed hybrid update and chart-labelled operator semantics}
This claim earns its public place by making a repair condition visible. The claim addresses this ambiguity: OC operators are typed update semantics; chart-labelled flow-one notation is admitted only for declared chart records, while proof/rewrite and guard/reset updates remain first-class non-smooth cases. No differentiability or ODE-solution theorem is promoted in current formal profile. The point of including it in the 1.3.3 release is to make the ambiguity auditable rather than to decorate the theory with another label.

The support class is Lean subset, structured proof sheet, and finite witness. The proof route is public proof sheet for the hybrid semantics theorem; the formal route is Lean declaration integrated\_operator\_semantics. In the public reading, this means that the claim is carried by a named theorem boundary, not by the fact that a row exists. If the proof sheet, Lean reference, finite witness, or mutation control does not carry the visible wording, the visible wording is the part that must move.

The claim boundary is Bounded to stated theorem assumptions, finite witnesses, and public falsifier boundary.. This boundary is not decorative. It prevents a bounded model-core statement from becoming a claim about unrestricted totality, full-scope prediction, or global priority. The scientific value of the claim is that it is strong enough to test a real structural distinction and narrow enough that an adversarial reviewer can name the condition that would reopen it.

A worked reading asks three local questions. Which typed field, operator, boundary, status predicate, or morphism is doing the work? Which situation is the positive case? Which negative case would have collapsed the distinction if the claim were only a slogan? The answer to those questions is the practical content of this subsection.

For T133-HYBRID, the corresponding falsifier is concrete rather than slogan-like: find a model satisfying the stated assumptions in which the asserted distinction collapses, show that the positive witness was accepted for the wrong reason, show that the negative control is non-responsive, or show that the public wording requires an assumption absent from the proof sheet.

The prior-art and empirical boundaries are local to the claim. Earlier formalisms may already carry part of the distinction; when they do, OC claims integration only for that part. Numeric replay rows can demonstrate artifact discipline, comparator behavior, and falsifier plumbing, but they do not by themselves complete a domain science. The reader should therefore leave with four handles: the conceptual problem, the typed construction, the evidence route, and the reopening route.

\subsection{Live-status cycle-mode requirement}
The reviewer-facing issue is whether the construction changes the space of admissible explanations. The claim addresses this ambiguity: Declared eligible-live status requires an explicit cycle mode or non-vacuous maintenance predicate. The point of including it in the 1.3.3 release is to make the ambiguity auditable rather than to decorate the theory with another label.

The support class is Lean subset, structured proof sheet, and finite witness. The proof route is public proof sheet for the cycle-mode theorem; the formal route is Lean declaration eligible\_live\_requires\_cycle\_or\_maintenance. In the public reading, this means that the claim is carried by a named theorem boundary, not by the fact that a row exists. If the proof sheet, Lean reference, finite witness, or mutation control does not carry the visible wording, the visible wording is the part that must move.

The claim boundary is Bounded to stated theorem assumptions, finite witnesses, and public falsifier boundary.. This boundary is not decorative. It prevents a bounded model-core statement from becoming a claim about unrestricted totality, full-scope prediction, or global priority. The scientific value of the claim is that it is strong enough to test a real structural distinction and narrow enough that an adversarial reviewer can name the condition that would reopen it.

A worked reading asks three local questions. Which typed field, operator, boundary, status predicate, or morphism is doing the work? Which situation is the positive case? Which negative case would have collapsed the distinction if the claim were only a slogan? The answer to those questions is the practical content of this subsection.

For T133-CYCLE, the corresponding falsifier is concrete rather than slogan-like: find a model satisfying the stated assumptions in which the asserted distinction collapses, show that the positive witness was accepted for the wrong reason, show that the negative control is non-responsive, or show that the public wording requires an assumption absent from the proof sheet.

The prior-art and empirical boundaries are local to the claim. Earlier formalisms may already carry part of the distinction; when they do, OC claims integration only for that part. Numeric replay rows can demonstrate artifact discipline, comparator behavior, and falsifier plumbing, but they do not by themselves complete a domain science. The reader should therefore leave with four handles: the conceptual problem, the typed construction, the evidence route, and the reopening route.

\subsection{Identity, residue, and rebirth classification}
The practical reading starts with the error that the claim prevents. The claim addresses this ambiguity: Identity continuation requires endpoint-bound identity evidence: identity class, declared invariant preservation, no residue token, equal source/target endpoint evidence, lifecycle identity-invariant truth, and typed source/target binding. Residue and rebirth evidence classes do not become identity continuation merely by preserving some invariants. No categorical Hom/composition theorem is promoted. The point of including it in the 1.3.3 release is to make the ambiguity auditable rather than to decorate the theory with another label.

The support class is Lean subset, structured proof sheet, and finite witness. The proof route is public proof sheet for the identity and rebirth theorem; the formal route is Lean declaration endpoint\_bound\_identity\_classification. In the public reading, this means that the claim is carried by a named theorem boundary, not by the fact that a row exists. If the proof sheet, Lean reference, finite witness, or mutation control does not carry the visible wording, the visible wording is the part that must move.

The claim boundary is Bounded to stated theorem assumptions, finite witnesses, and public falsifier boundary.. This boundary is not decorative. It prevents a bounded model-core statement from becoming a claim about unrestricted totality, full-scope prediction, or global priority. The scientific value of the claim is that it is strong enough to test a real structural distinction and narrow enough that an adversarial reviewer can name the condition that would reopen it.

A worked reading asks three local questions. Which typed field, operator, boundary, status predicate, or morphism is doing the work? Which situation is the positive case? Which negative case would have collapsed the distinction if the claim were only a slogan? The answer to those questions is the practical content of this subsection.

For T133-ID, the corresponding falsifier is concrete rather than slogan-like: find a model satisfying the stated assumptions in which the asserted distinction collapses, show that the positive witness was accepted for the wrong reason, show that the negative control is non-responsive, or show that the public wording requires an assumption absent from the proof sheet.

The prior-art and empirical boundaries are local to the claim. Earlier formalisms may already carry part of the distinction; when they do, OC claims integration only for that part. Numeric replay rows can demonstrate artifact discipline, comparator behavior, and falsifier plumbing, but they do not by themselves complete a domain science. The reader should therefore leave with four handles: the conceptual problem, the typed construction, the evidence route, and the reopening route.

\subsection{Declared component independence of the semantic verdict interface}
The first question for this claim is what weaker vocabulary would confuse. The claim addresses this ambiguity: Within the declared current release tuple semantics, each tuple component has a one-field semantic keep/drop witness that changes the release verdict. The point of including it in the 1.3.3 release is to make the ambiguity auditable rather than to decorate the theory with another label.

The support class is Lean subset, structured proof sheet, and finite witness. The proof route is public proof sheet for the minimality witness theorem; the formal route is Lean declaration release\_tuple\_semantic\_component\_irredundant. In the public reading, this means that the claim is carried by a named theorem boundary, not by the fact that a row exists. If the proof sheet, Lean reference, finite witness, or mutation control does not carry the visible wording, the visible wording is the part that must move.

The claim boundary is Bounded to stated theorem assumptions, finite witnesses, and public falsifier boundary.. This boundary is not decorative. It prevents a bounded model-core statement from becoming a claim about unrestricted totality, full-scope prediction, or global priority. The scientific value of the claim is that it is strong enough to test a real structural distinction and narrow enough that an adversarial reviewer can name the condition that would reopen it.

A worked reading asks three local questions. Which typed field, operator, boundary, status predicate, or morphism is doing the work? Which situation is the positive case? Which negative case would have collapsed the distinction if the claim were only a slogan? The answer to those questions is the practical content of this subsection.

For T133-MIN, the corresponding falsifier is concrete rather than slogan-like: find a model satisfying the stated assumptions in which the asserted distinction collapses, show that the positive witness was accepted for the wrong reason, show that the negative control is non-responsive, or show that the public wording requires an assumption absent from the proof sheet.

The prior-art and empirical boundaries are local to the claim. Earlier formalisms may already carry part of the distinction; when they do, OC claims integration only for that part. Numeric replay rows can demonstrate artifact discipline, comparator behavior, and falsifier plumbing, but they do not by themselves complete a domain science. The reader should therefore leave with four handles: the conceptual problem, the typed construction, the evidence route, and the reopening route.

\subsection{Adjacent K-level witness consistency}
The claim is useful only if it separates two situations that a less typed account would merge. The claim addresses this ambiguity: Every declared adjacent K-level transition K0-\textgreater{}K12 has a release-atlas row, retained-witness evaluator check, executable finite row, and inert-witness demotion control inside the current formal profile release classifier; independent semantic irreducibility beyond this declared classifier is a future proof obligation, not a promoted current formal profile theorem. The point of including it in the 1.3.3 release is to make the ambiguity auditable rather than to decorate the theory with another label.

The support class is Lean subset, structured proof sheet, and finite witness. The proof route is public proof sheet for the K-level witness theorem; the formal route is Lean declaration release\_atlas\_manifest\_has\_total\_finite\_case\_coverage. In the public reading, this means that the claim is carried by a named theorem boundary, not by the fact that a row exists. If the proof sheet, Lean reference, finite witness, or mutation control does not carry the visible wording, the visible wording is the part that must move.

The claim boundary is Bounded to stated theorem assumptions, finite witnesses, and public falsifier boundary.. This boundary is not decorative. It prevents a bounded model-core statement from becoming a claim about unrestricted totality, full-scope prediction, or global priority. The scientific value of the claim is that it is strong enough to test a real structural distinction and narrow enough that an adversarial reviewer can name the condition that would reopen it.

A worked reading asks three local questions. Which typed field, operator, boundary, status predicate, or morphism is doing the work? Which situation is the positive case? Which negative case would have collapsed the distinction if the claim were only a slogan? The answer to those questions is the practical content of this subsection.

For T133-KLEVEL, the corresponding falsifier is concrete rather than slogan-like: find a model satisfying the stated assumptions in which the asserted distinction collapses, show that the positive witness was accepted for the wrong reason, show that the negative control is non-responsive, or show that the public wording requires an assumption absent from the proof sheet.

The prior-art and empirical boundaries are local to the claim. Earlier formalisms may already carry part of the distinction; when they do, OC claims integration only for that part. Numeric replay rows can demonstrate artifact discipline, comparator behavior, and falsifier plumbing, but they do not by themselves complete a domain science. The reader should therefore leave with four handles: the conceptual problem, the typed construction, the evidence route, and the reopening route.


\section{Scientific Reading Protocol for the 1.3.3 Monolith}
\label{sec:oc133-scientific-reading-protocol}
The monolith is intentionally long because it has to serve several readers at once: a formal reviewer, a domain scientist, an editor, a reproducibility auditor, and a hostile reader looking for overclaim. The recommended protocol is not to read every evidence row first. The reader should first understand the object grammar, then the theorem ceiling, then the executable evidence, then the comparator boundary, and only then the release and journal-preparation apparatus.

Step one is the object grammar. The reader should be able to say what a carrier is, what a realization is, what lawful possibility permits, what liveness means at a time slice, what residue preserves after death, what an identity morphism is allowed to preserve, and why a boundary is not merely a metaphor. If this grammar is unclear, later empirical and release sections should be paused rather than skimmed.

Step two is the theorem ceiling. Each theorem should be read as a bounded claim with assumptions, proof route, finite witness route, and counterexample boundary. The purpose of the theorem registry is not to impress the reader with labels; it is to make it impossible for a public sentence to float away from its proof obligation. A label that cannot be traced to proof and boundary is editorially unfinished.

Step three is the executable layer. The Lean subset and finite semantic checks are not decorations and they are not substitutes for the entire mathematical theory. They are selected executable anchors: places where the reader can see that a semantic distinction has been encoded, accepted, rejected, or mutation-tested. Their scientific role is strongest when the negative control is as visible as the positive witness.

Step four is the empirical layer. The target-blind rows and replay-QA tables are read as evidence of disciplined claim-to-data plumbing. They demonstrate formula binding, snapshot binding, comparator binding, uncertainty or residual accounting, negative control, falsifier, and replay hash. They do not close every future domain; they show how a domain claim must be made if it is to enter the release surface.

Step five is the comparator layer. The reader should look for accepted overlap before looking for residual delta. If the manuscript sounds as if it invented general systems theory, autopoiesis, dynamical systems, category theory, RAF theory, complexity measures, identity theory, or reproducibility engineering, then the wording has failed. The defensible contribution is the bounded integration and governance of claim, proof, data, comparator, and reopening conditions.

Step six is the adversarial layer. Reviewer objections are not an appendix of public relations. They are part of the scientific control surface. A good objection names the attacked claim, the exact criticism route, the evidence that would answer it, and the residual condition that would reopen it. The release is stronger when those routes are explicit because a hostile reader can test the same paths the authors used.

Step seven is the publication layer. GitHub and Zenodo package the public scientific object; journal packets remain owner-review preparation material until a later venue-specific action. Metadata, checksums, PDFs, and asset lists are not secondary chores: they are part of the reproducibility perimeter. A release whose metadata contradicts its manuscript is scientifically damaged even if the proofs are locally sound.

This reading protocol is also a repair protocol. When a defect appears, the correction should target the first layer where the defect originates: ontology, theorem boundary, executable witness, empirical replay, comparator positioning, adversarial response, editorial prose, or publication packaging. That is how the system avoids random review churn and closes the most important vulnerabilities before spending attention on polish.
